\section{第三日第三讲}

本讲第一个目标，是对一个球 $G$-簇 $X$ 联系一个映射
\begin{equation}
  \mathrm{SL}_2\times\check G_X\to\check G.
  \label{eq:10-1}
\end{equation}
这个映射的动机来自对 $H_G$-模 $C_c^\infty(X(F))^{G(O)}$ 的研究：对「好」的 $X$，我们预期有同构
\begin{equation}
  C_c^\infty(X(F))^{G(O)}\cong\mathbb C[\Rep\check G_X].
  \label{eq:10-2}
\end{equation}
它是 $H_G\simeq\mathbb C[\Rep\check G]$ 上的模同构。映射 (10.1) 用来描述这个同构的右端。在解释右端之前，我们先通过定义 $X$ 的 Cartan 子群 $A_X$ 与 Weyl 群 $W_X$，来定义与 $X$ 相联系的嵌入 $\check G_X\hookrightarrow\check G$。

\subsection{相对 Cartan 子群}

设 $B$ 是 $G$ 的 Borel 子群，$X^\circ$ 是相应的开 $B$-轨道。

对 $G$ 可以定义一个环面 $A$，称为 $G$ 的\textbf{万有 Cartan 子群}，使得对任意 Borel 子群 $B$ 连同其幂幺根 $N$，都有典范同构 $B/N\xrightarrow{\sim}A$。由于 $B$ 作用在 $X^\circ$ 上，$A$ 作用在 $X^\circ/N$ 上，而这个作用经过 $A$ 的一个商 $A_X$ 分解。

$A_X$ 的特征群 $X^*(A_X)$ 可以与 $B$-模 $k(X)$ 中出现的 $B$-本征特征标组成的群等同起来。事实上，由于 $X$ 有开稠的 $B$-轨道，$k(X)$ 的 $B$-本征特征标 $\chi\in X^*(B)$ 由相应的非零 $B$-本征函数 $f_\chi\in k(X)^{(B)}$ 决定，而后者在 $k$ 的意义下只差一个标量（即短正合列 (10.3)）。我们还有
\[
  k(X)^{(B)}=k(X^\circ)^{(B)}=k(B_x\backslash B)^{(B)}=k(B_x\backslash B/N)^{(B/N)}=k(X^\circ/N)^{(A)},
\]
其中 $B$-轨道 $X^\circ\simeq B_x\backslash B$，$x\in X^\circ$，$B_x$ 是相应的稳定化子。

在把 $X^*(A_X)$ 等同于 $k(X)$ 的 $B$-本征特征标群之下，我们得到短正合列
\begin{equation}
  1\to k^*\to k(X)^{(B)}\xrightarrow[\ f_\chi\mapsto\chi\ ]{}X^*(A_X)\to 1.
  \label{eq:10-3}
\end{equation}

\subsection{小 Weyl 群}

对一个球 $G$-簇 $X$，小 Weyl 群 $W_X$ 有三种等价的定义/看法。最后一种最便于计算。

\subsubsection{来自 $G$-不变赋值}

设 $V$ 是 $k(X)$ 上 $G$-不变的 $\mathbb Q$-值离散赋值的集合。利用正合列 (10.3)，Knop 证明了有单射
\[
  V\hookrightarrow\Hom(X^*(A_X),\mathbb Q)=X_*(A_X)\otimes\mathbb Q=:\mathfrak a_X,\qquad v\mapsto(f_\chi\mapsto v(f_\chi)).
\]

\begin{prop}
[Knop，见《The Luna--Vust theory of spherical embeddings》]
存在子群 $W_X\subset W$，它稳定 $X^*(A_X)\subset X^*(A)$，并且 $V$ 是 $W_X$ 在 $\mathfrak a_X$ 上作用的基本域。
\end{prop}

\subsubsection{来自余切丛的 GIT 商}

\begin{prop}
[Knop]
设 $X$ 是齐性、拟仿射的球 $G$-簇，$\mathfrak a_X=\mathrm{Lie}\,A_X$，$\mathfrak a=\mathrm{Lie}\,A$，于是 $\mathfrak a_X^*\subset\mathfrak a^*$。存在稳定 $\mathfrak a_X^*$ 的子群 $W_X\subset W$，以及典范同构
\[
  T^*X\gquot G\xrightarrow{\ \sim\ } \mathfrak a_X^*\gquot W_X.
\]
此外，映射 $T^*X\to\mathfrak g^*\to\mathfrak a^*\gquot W$ 经过上述映射分解。
\end{prop}

\begin{exmp}
[群情形]
设 $X=H$，$G=H\times H$ 按 $h\cdot(h_1,h_2)=h_1^{-1}hh_2$ 右作用。则
\[
  T^*H\gquot G=T^*_eH\gquot\mathrm{Stab}_G(e)=\mathfrak h^*\gquot H,
\]
而由 Chevalley 同构 $\mathfrak h^*\gquot H=\mathfrak a_H^*\gquot W_H$，其中 $\mathfrak a_H$ 是 $H$ 中极大环面的 Lie 代数，$W_H$ 是 $H$ 的 Weyl 群。于是 $W_H$ 就是 $X$ 的小 Weyl 群。
\end{exmp}

\subsubsection{来自 Knop 在 Borel 轨道上的 Weyl 群作用}

Knop 在他的论文「On the set of orbits for a Borel subgroup」中定义了 $W_G$ 在 $X$ 的 $B$-轨道上的如下作用：

\begin{defn}
[$W_G$ 在 $X$ 的 $B$-轨道上的作用]
设 $\alpha\in\Delta_G$ 是 $G$ 的单根，$w_\alpha\in W_G$ 是相应的反射。设 $\gamma$ 是 $X$ 中的 $B$-轨道，$P_\alpha\supset B$ 是 $\alpha$ 对应的极小抛物子群，则 $w_\alpha\cdot\gamma$ 将是 $\gamma P_\alpha$ 中的 $B$-轨道，定义如下：$\gamma P_\alpha/R(P_\alpha)$ 是齐性球 $P_\alpha/R(P_\alpha)\simeq\mathrm{PGL}_2$-簇。齐性球 $\mathrm{PGL}_2$-簇在等价意义下只有 $4$ 种类型，我们用下表在每种类型的 Borel 轨道上定义 $w_\alpha$ 的作用：

\smallskip
\noindent\begin{tabular}{|c|c|p{0.30\textwidth}|}
  \hline
  类型 & 齐性球 $\mathrm{PGL}_2$-簇 & Borel 轨道与 $w_\alpha$ 的作用\\
  \hline
  $G$ & $\mathrm{pt}$ & 一个闭轨道\\
  \hline
  $T$ & $\mathbb G_m\backslash\mathrm{PGL}_2$ & 两个闭轨道、一个开轨道；$w_\alpha$ 交换两个闭轨道、不动开轨道\\
  \hline
  $N$ & $N(\mathbb G_m)\backslash\mathrm{PGL}_2$ & 一个闭轨道、一个开轨道；$w_\alpha$ 不动这两个轨道\\
  \hline
  $U$ & $N\cdot S\backslash\mathrm{PGL}_2$，其中 $N=\begin{pmatrix}1&a\\0&1\end{pmatrix}\subset S\subset B$ & 一个开轨道、一个闭轨道；$w_\alpha$ 交换这两个轨道\\
  \hline
\end{tabular}
\smallskip

合起来，这就定义了 $G$ 的 Weyl 群 $W_G$ 在 $X$ 的 $B$-轨道上的一个作用。
% [Tutor Note: 原文 U 行写作 "one open and open closed orbit"，词序有误；按上下文（$w_\alpha$ 交换两个轨道）应为「一个开轨道、一个闭轨道」，译文按此理解，已在 reports/ch10.md 登记。]
\end{defn}

\begin{defn}
给定球 $G$-簇 $X$，可以联系一个抛物子群 $P(X):=\{g\in G\mid X^\circ g=X^\circ\}$ 连同其 Levi 子群 $L(X)$。若 $P(X)=B$，则称 $X$ 是\textbf{tempered} 的。
\end{defn}

\begin{exmp}
[tempered 性]
群情形是 tempered 的。
\end{exmp}

小 Weyl 群在 Knop 的作用之下这样出现：

\begin{prop}
[Knop]
在此作用下，开 $B$-轨道 $X^\circ$ 的稳定化子是
\[
  W_X\ltimes W_{P(X)},
\]
其中 $W_{P(X)}$ 是 $P(X)$ 的 Weyl 群。
\end{prop}

我们将用这个命题来计算某些 $X$ 的小 Weyl 群。

\begin{exmp}
[群情形]
设 $X=H$，$G=H\times H$ 按 $h\cdot(h_1,h_2)=h_1^{-1}hh_2$ 右作用。另一种看法是把 $X$ 看作 $X=H_{\mathrm{diag}}\backslash H\times H$，而 $G=H\times H$ 按右乘作用。

由 Bruhat 分解，$X$ 的 $B\times B$-轨道是 $BwB$（$w\in W_H$），其中 $W_H$ 是 $H$ 的 Weyl 群。开 $B\times B$-轨道是 $Bw_0B$，其中 $w_0\in W_H$ 是由 $B\subset H$ 定义的最长元素。

对 $H\times H$ 的单根 $(\alpha,0)\in\Delta_{H\times H}=\Delta_H\sqcup\Delta_H$，相应的极小抛物子群是 $P_\alpha\times B$。考虑
\[
  (Bw_0B)\cdot(P_\alpha\times B)/R(P_\alpha\times B)=R(P_\alpha)\backslash P_\alpha w_0B/B.
\]
这总是类型 $U$\footnote{可以证明 $R(P_\alpha)w_0B\in R(P_\alpha)\backslash P_\alpha w_0B/B$ 在 $R(P_\alpha)\backslash P_\alpha$ 作用下的稳定化子是 $R(P_\alpha)\backslash B^-\cap P_\alpha$，因此有同构
  $\mathrm{PGL}_2/B^-\xrightarrow{\sim}R(P_\alpha)\backslash P_\alpha/(B^-\cap P_\alpha)\xrightarrow{[g]\mapsto[gw_0]}R(P_\alpha)\backslash P_\alpha w_0B/B$。}（见上表）。特别地，我们知道
$w_{(\alpha,0)}=(w_\alpha,1)\in W_{H\times H}=W_H\times W_H$ 在 $P_\alpha w_0B$ 中交换哪两个 Borel 轨道。事实上，由于 $P_\alpha=B\sqcup Bw_\alpha B$ 且 $w_0(\alpha)\in\Phi^-$，故\footnote{见 Tits 系统下 Bruhat 分解的证明。}
\[
  P_\alpha w_0B=Bw_0B\sqcup Bw_\alpha w_0B.
\]
类似地，$w_{(0,\alpha)}=(1,w_\alpha)\in W_{H\times H}=W_H\times W_H$ 交换 $Bw_0B$ 与 $Bw_0w_\alpha B=Bw_\alpha w_0B$（因为对任意 $w\in W$、$\alpha\in\Phi$ 有 $ww_\alpha w^{-1}=w_{w(\alpha)}$）。这说明 $(w_\alpha,w_\alpha)$ 不动开 Borel 轨道 $X^\circ=Bw_0B$。此外，$(w,w')\in W_{H\times H}$ 不动 $X^\circ$ 当且仅当 $w=w'$。由于 $X$ 是 tempered 的，$W_{P(X)}=1$，于是由 Knop 的作用得到 $W_X=\{(w,w):w\in W_H\}$。
\end{exmp}

\begin{exmp}
[$X=\mathrm{PGL}_2^{\mathrm{diag}}\backslash\mathrm{PGL}_2^3$ 作为 $\mathrm{PGL}_2^3$-空间]
记 $B$ 为 $\mathrm{PGL}_2$ 中的上三角矩阵。先计算 $X$ 的开 Borel 轨道，即 $\mathrm{PGL}_2^{\mathrm{diag}}\backslash\mathrm{PGL}_2^3/B^3$ 中最大的双陪集。把 $\mathrm{PGL}_2^3/B^3$ 与 $(\mathbb P^1)^3$ 等同，容易看出最大轨道对应于代表元 $x=([1,0],[0,1],[1,1])\in(\mathbb P^1)^3$：因为 $\mathrm{PGL}_2$ 在 $\mathbb P^1$ 上可迁地作用，要得到大轨道就需要三个互不相同的点。

先通过计算稳定化子 $B_x$ 来求 $A_X$。等同 $\mathrm{PGL}_2/B$ 与 $\mathbb P^1$ 由
\[
  \begin{pmatrix}a&*\\c&*\end{pmatrix}\leftrightarrow[a,c]
\]
给出。于是经过一些矩阵乘法（注意这里 $B^3$ 从右作用），可以看出 $B_x=\{1\}$，从而 $A_X=\mathbb G_m^3$。

现在计算 $W_X$。对 $(1,1,w)$（$w$ 是 $W_{\mathrm{PGL}_2}$ 中的非平凡元素），相应的极小抛物子群是 $B\times B\times\mathrm{PGL}_2$。取 $x$ 在 $\mathrm{PGL}_2^3$ 中的一个代表元，计算商
\[
  \begin{aligned}
    Y_2&=\mathrm{PGL}_2^{\mathrm{diag}}\Big\backslash\Big(\begin{pmatrix}1&0\\0&1\end{pmatrix}B,\ \begin{pmatrix}0&1\\1&0\end{pmatrix}B,\ \begin{pmatrix}1&0\\1&1\end{pmatrix}B\Big)\cdot B\times B\times\mathrm{PGL}_2\,/\,B\times B\times\{1\}\\
    &=\mathrm{PGL}_2^{\mathrm{diag}}\backslash B\times wB\times\mathrm{PGL}_2\,/\,B\times B\times\{1\}\\
    &\cong\mathrm{PGL}_2\backslash\mathrm{PGL}_2=\text{点}.
  \end{aligned}
\]
因此 $(1,1,w)$ 稳定开 Borel 轨道。另两个单根的情形类似。可以验证与 $X$ 相联系的抛物子群是 $B^3$，从而 $W_X=W_{\mathrm{PGL}_2^3}$。
\end{exmp}

\begin{exmp}
[练习中的其他例子]\leavevmode\par
\begin{enumerate}[label=(\arabic*),leftmargin=1.8em,itemsep=0.2em]
  \item $X=\mathrm{Mat}_n$ 作为 $\mathrm{GL}_n\times\mathrm{GL}_n$-空间。
  \item $X=\mathrm{SO}_n\backslash\mathrm{SO}_{n+1}$ 作为 $\mathrm{SO}_{n+1}$-空间。
  \item $X=\mathrm{SO}_{n+1}$ 作为 $\mathrm{SO}_n\times\mathrm{SO}_{n+1}$-空间。预期答案：$\check G_X=\check G$。
\end{enumerate}
\end{exmp}

\subsection{球簇的对偶群}

对偶群 $\check G_X$ 看作 $\check G$ 的子群，由它的 Cartan 子群 $A_X$ 与 Weyl 群 $W_X$ 完全确定。

\begin{exmp}
[群情形]
当 $X=H$、$G=H\times H$ 时，可以证明 $\check G_X=\check H$。
\end{exmp}

\begin{exmp}
当 $X=\mathrm{PGL}_2^{\mathrm{diag}}\backslash\mathrm{PGL}_2^3$ 时，$\check G_X=\check G$。
\end{exmp}

\begin{exmp}
[Hecke]
当 $X=\mathbb G_m\backslash\mathrm{PGL}_2$ 时……
\end{exmp}

\subsection{$\mathrm{SL}_2$-项}

见 BZSV，§4。

我们已经定义了 $i:\check G_X\to\check G$。由此 $\mathbb C[\Rep\check G_X]$ 是 $\mathbb C[\Rep\check G]$ 上的模，其中 $V\in\Rep\check G$ 通过 $i^*V\otimes W$ 作用在 $W\in\Rep\check G_X$ 上。然而这并非给出 (10.2) 的正确模结构。需要对 $i$ 作一个移位，下面的例子说明这一点。

\begin{exmp}
[Whittaker]
设 $X=\mathrm{pt}=G\backslash G$。$H_G$ 通过平凡特征标 $\chi_1:H_G\to\mathbb C$ 作用在 $C_c^\infty(X(F))^{G(O)}=\mathbb C$ 上。在 Satake 同构 $\Hom(H_G,\mathbb C)=(\check A\gquot W)(\mathbb C)$ 之下，$\chi_1$ 对应于 $q^\rho\in\check A(\mathbb C)$。

如何解释这个 $q^\rho$：$\rho\in\check{\mathfrak g}$ 对应于 $\rho=d\tilde\rho$，其中 $\tilde\rho:\mathbb G_m\to\check G$，而 $q^\rho=\tilde\rho(q)$。

可以对 $X$ 联系一个 $\mathfrak{sl}_2$-三元组 $(h,e,f)$，它在 $\check L(X)\subset\check G$（$P(X)\subset G$ 对应的 Levi 子群）中是主三元组。这个三元组与 $\check G_X$ 交换，从而定义 $\iota:q^{h/2}\check G_X\hookrightarrow\check G$，它提升为 $\mathrm{SL}_2\times\check G_X\to\check G$。$i$ 的这个移位 $\iota$ 给出 (10.2)。
\end{exmp}

\begin{exmp}
[Whittaker]
验证这个移位在作用的意义下是正确的。$\mathfrak{sl}_2$-三元组 $(h,e,f)$ 是主三元组，其中 $P(X)=\cdots$
\end{exmp}

\subsection{超球簇}

本讲第二个目标是推广球簇。由此引出超球簇的概念。基本的例子是球簇的余切空间。

\begin{defn}
带交换 $G_{\mathrm{gr}}$ 的 Hamilton $G$-簇 $M$ 称为\textbf{超球}的，若它满足下列条件：
\begin{enumerate}[label=(\arabic*),leftmargin=1.8em,itemsep=0.15em]
  \item $M$ 光滑且仿射。
  \item 作为 $G$-簇，$M$ 是余迷向的。
  \item $G$ 在 $M$ 上作用的通常稳定化子连通。
  \item 矩映射 $M\to\mathfrak g^*$ 的像与幂零锥相交（即交非空）。
  \item 分次作用是「中性的」。特别地，它与 $\mathfrak g^*$ 上的平方作用相容。
\end{enumerate}
\end{defn}

\begin{rem}\leavevmode\par
\begin{itemize}[leftmargin=1.4em]
  \item 关于 (4)：若 $M=T^*X$，则矩映射的像总包含零截面。
  \item 实验表明 (2)+(3) 与 (1) 对偶。
\end{itemize}
\end{rem}

\begin{exmp}
[$M$ 不是 $T^*X$ 形状的情形]
考虑 $T^*(N\backslash G)$，其中 $N$ 是 $B$ 的幂幺根。我们用 Hamilton 约化来解释商空间的余切空间，其定义为矩映射下 $\{0\}$ 的纤维对 $N$ 的商：
\[
  T^*(N\backslash G)\overset{\text{def}}{=}N\gquot G=\{0\}\times_N\mathfrak n^*T^*G.
\]
也就是说，我们看的是 $G$ 上与 $N$-轨道正交的余切向量。注意这个空间不是仿射的。

略作修改即得 Whittaker 余切空间。构造它时，先取 $\mathfrak g\simeq\mathfrak g^*$ 中的 $\mathfrak{sl}_2$-三元组 $(h,e,f)$。我们要 $f$ 在 $\mathfrak g^*$ 中正则幂零，使 $f$ 给出 $\mathfrak n$ 的非退化特征标。令 Whittaker 余切空间为
\[
  T^*((N,f)\backslash G)=\{f\}\times_N\mathfrak n^*T^*G=M.
\]
这样 $M$ 是仿射的，且是超球的。注意它还带有 Kostant 截面
\[
  M\gquot G=\mathfrak a^*\gquot W\xrightarrow{\ \text{Kostant}\ }M.
\]
这给出 $M\simeq(\mathfrak a^*\gquot W)\times G$。
\end{exmp}

\begin{thm}
[超球簇的结构定理]
超球 Hamilton $G$-簇 $M$ 由下列数据搭成：
\[
  (G,H,\iota,S),
\]
其中 $G$ 是约化群，$H$ 是 $G$ 的约化子群，$\iota$ 是 $G$ 中的一个 $\mathfrak{sl}_2$-三元组，$S$ 是 $G$ 的辛表示。
\end{thm}

上述构造称为 Whittaker 诱导，记作 $M=\mathrm{W}_\iota\mathrm{Ind}_H^G(S)$。

对这样的空间 $M$，我们期待某个对偶空间 $\check M$。

当 $M=T^*X$（$X$ 球簇）时，我们有如下对应
\[
  \check M\ \longleftrightarrow\ (\check G,\check G_X,\iota_{P(X)},\check S),
\]
其中 $\check G$ 是对偶群，$\check G_X$ 是此前定义的与 $X$ 相联系的约化子群，$\iota_{P(X)}$ 是 $P(X)$ 对应的 $\mathfrak{sl}_2$-三元组，$\check S$ 是由 $X$ 的颜色（即非 $G$-稳定的 $B$-稳定除子）决定的辛 $\check G$-表示（在同构意义下）。非典范地有
\[
  \check M\simeq V_X\times_{\check G_X}\check G,\qquad V_X=\check S\oplus(\check{\mathfrak g}/\check{\mathfrak g}_X)_e+\check{\mathfrak g}_1,
\]
其中 $e$ 来自 $\mathfrak{sl}_2$-三元组 $\iota_{P(X)}$，而 $\check{\mathfrak g}_1$ 是 $\check{\mathfrak g}$ 中的 $1$ 次部分。

对任意 $\mathfrak{sl}_2$-三元组 $\iota:\mathfrak{sl}_2\hookrightarrow\mathfrak g$，$h$ 通过 $[h,-]$ 的本征值给出 $\mathfrak g$ 的分次
\[
  \mathfrak g=\mathfrak g_+\oplus\mathfrak g_0\oplus\mathfrak g_-,
\]
其中总有 $f\in\mathfrak g_{-2}$。若这个 $\mathfrak{sl}_2$-三元组给出偶分解（即只有偶次），则在等同 $\mathfrak g\simeq\mathfrak g^*$ 下，$f$ 定义一个特征标 $f:\mathfrak g_+\to F$。一般情形下，$f$ 定义一个特征标 $\mathfrak g_{\ge2}\times\mathfrak g_{\ge2}\to F$（Heisenberg 群）。
